\chapter{คลังชุดคำสั่งวิศวกรรมพรอมพต์สำหรับการพัฒนาแอปพลิเคชัน}
\label{app:prompts}

ภาคผนวกนี้รวบรวมคลังชุดคำสั่งวิศวกรรมพรอมพต์ (Prompt Engineering Corpus) ตามกรอบการทำงาน CLEAR Prompting Framework เพื่อให้ผู้เริ่มต้น นักศึกษา และนักพัฒนา สามารถนำไปใช้สั่งการโมเดลปัญญาประดิษฐ์เชิงกำเนิด (เช่น Claude 3.7 Sonnet, Gemini 1.5/2.0 Flash/Pro, หรือ GPT-4o) ในการสร้างโค้ดและแก้ไขปัญหาในแต่ละโมดูลของระบบ \textbf{SoilpH-HT AI} ได้อย่างถูกต้อง แม่นยำ และเป็นไปตามหลักการทางวิทยาศาสตร์

\section{ชุดพรอมพต์ออกแบบสถาปัตยกรรมระบบ Clean Architecture}

\begin{promptbox}[title={พรอมพต์ ก.1 การวางโครงสร้างโปรเจกต์ Flutter Clean Architecture}]
\begin{lstlisting}
คุณคือสถาปนิกซอฟต์แวร์อาวุโสด้าน Mobile Embedded System
จงออกแบบโครงสร้างโปรเจกต์ Flutter สำหรับแอปพลิเคชัน SoilpH-HT AI วิเคราะห์ค่ากรด-ด่างดิน
โดยใช้หลักการ Clean Architecture ที่แบ่งความรับผิดชอบอย่างชัดเจน
1. Core Layer ได้แก่ ค่าคงที่, ธีมสี Cyber Dark Glow, เกณฑ์วิเคราะห์ความเป็นกรด-ด่างและใบสั่งปูนโดโลไมต์
2. Data Layer ได้แก่
   - บริการสื่อสารพอร์ต USB Type-C OTG Serial (CH340/MAX485)
   - โมดูลถอดรหัสเฟรมข้อมูล Modbus RTU สำหรับ pH อุณหภูมิ และความชื้น พร้อมตรวจสอบ CRC-16
   - บริการบันทึกภาพถ่ายดินพร้อมพิกัด GPS และเมทาดาทา JSON
3. Domain Layer ได้แก่
   - เอนจินโมเดลโครงข่ายประสาทเทียมอิงฟิสิกส์ (Physics-Informed Neural Network หรือ PINN)
   - โมเดลฟิสิกส์การปรับเทียบอุณหภูมิตามสมการเนิร์นสต์ (Nernstian pH) และอิมพีแดนซ์ดินแห้ง
   - เอนจินวิเคราะห์สุขภาพดินและคำนวณอัตราปูนโดโลไมต์สำหรับสวนทุเรียน
4. Presentation (UI) Layer ได้แก่
   - แดชบอร์ดแสดงผลมาตรวัด pH แบบวงแหวนเรืองแสง PhCircularGauge
   - วิดเจ็ตแบบ Responsive ปรับขนาดอัตโนมัติ ไม่ล้นจอ
   - ระบบกล้อง AI Vision พร้อมแผงเทเลเมทรี HUD และพิกัดดาวเทียม
จงแสดงโครงสร้างไดเรกทอรีไฟล์ทั้งหมดในรูปแบบ Tree พร้อมคำอธิบายหน้าที่ของแต่ละไฟล์
\end{lstlisting}
\end{promptbox}

\section{ชุดพรอมพต์สร้างไดรเวอร์ USB OTG และถอดรหัส Modbus RTU}

\begin{promptbox}[title={พรอมพต์ ก.2 การเขียนโมดูลถอดรหัส Modbus RTU และตรวจสอบ CRC-16}]
\begin{lstlisting}
คุณคือวิศวกรระบบสมองกลฝังตัวผู้เชี่ยวชาญโปรโตคอลอุตสาหกรรม
จงเขียนคลาส ModbusParser และ SoilPhReading ภาษา Dart สำหรับถอดรหัสข้อมูลจากหัววัด Soil Parameter Tester
1. รองรับโครงสร้างเฟรมมาตรฐาน Modbus RTU ได้แก่
   [Slave ID (1B), Function Code (1B), Byte Count (1B), Data (8B), CRC-16 (2B)]
2. เขียนอัลกอริทึมคำนวณและตรวจสอบ CRC-16 Modbus (พหุนาม 0xA001) แบบบิตต่อบิต
3. ใช้ ByteData เพื่ออ่านค่ารีจิสเตอร์ขนาด 16-bit Unsigned Integer (Big-Endian)
4. แปลงค่าตามตัวคูณมาตราส่วน (Scale Factors) ของหัววัด ได้แก่
   - อุณหภูมิดิน หารด้วย 10.0 (องศาเซลเซียส)
   - ความชื้นดิน หารด้วย 10.0 (เปอร์เซ็นต์)
   - สภาพนำไฟฟ้า ค่าจำนวนเต็ม (ไมโครซีเมนต์ต่อเซนติเมตร)
   - ความเป็นกรด-ด่าง หารด้วย 10.0 (pH)
5. หากตรวจสอบ CRC-16 ไม่ผ่าน ให้โยนข้อยกเว้น FormatException พร้อมแสดง Hex Dump
\end{lstlisting}
\end{promptbox}

\section{ชุดพรอมพต์สร้างโมเดลโครงข่ายประสาทเทียมอิงฟิสิกส์ (PINN)}

\begin{promptbox}[title={พรอมพต์ ก.3 การสร้างโมเดล Two-Stage Decoupling PINN บนอุปกรณ์ฝังตัว}]
\begin{lstlisting}
คุณคือนักวิทยาศาสตร์ข้อมูลด้านฟิสิกส์เกษตรดิจิทัล (Digital Agriphysics Data Scientist)
จงเขียนคลาส DeepLearningCalibrator ภาษา Dart เพื่อชดเชยค่าความคลาดเคลื่อนของกรด-ด่างดิน ดังนี้
1. การชดเชยขั้นตอนที่ 1 (First-Principles Physics Stage)
   - ชดเชยความชันเนิร์นสเตียนตามอุณหภูมิ S_N(T) = 2.3026 * R * T / F
   - ชดเชยการเลื่อนของการแตกตัวของน้ำ (pKw drift): 0.017 * (25.0 - T)
   - ชดเชยอิมพีแดนซ์รอยต่อในสภาวะดินแห้ง 0.35 * exp(-0.08 * Moisture)
2. การชดเชยขั้นตอนที่ 2 (Residual MLP PINN Stage)
   - โครงข่ายประสาทเทียม MLP 3 ชั้น ชดเชยปฏิสัมพันธ์ร่วมไม่เป็นเชิงเส้นระหว่างอุณหภูมิและความชื้น
   - ใช้ฟังก์ชันกระตุ้น Swish และ GELU เพื่อความราบเรียบเชิงอนุพันธ์
   - ควบคุมขอบเขตผลลัพธ์ด้วยคำสั่ง clamp(3.0, 10.5) เพื่อป้องกันค่าหลุดออกนอกช่วงทางกายภาพ
3. ฟังก์ชันต้องคืนค่า Map<String, double> ประกอบด้วยค่าดิบ ค่าปรับเทียบ และค่าส่วนต่าง (Delta)
\end{lstlisting}
\end{promptbox}

\section{ชุดพรอมพต์ออกแบบหน้าจอแดชบอร์ดและวิดเจ็ตมาตรวัด pH แบบ Responsive}

\begin{promptbox}[title={พรอมพต์ ก.4 การสร้างวิดเจ็ตมาตรวัด pH ทรงกลม CustomPainter แบบ Responsive}]
\begin{lstlisting}
คุณคือผู้เชี่ยวชาญการออกแบบ Flutter UI/UX สไตล์ Glassmorphism และ Responsive Design
จงสร้างคลาส PhCircularGauge และ CustomPainter สำหรับแสดงผลค่าความเป็นกรด-ด่างของดิน ดังนี้
1. มีวงแหวนความคืบหน้า (Circular Arc) 240 องศา แสดงระดับค่า pH จาก 3.0 ถึง 10.0
2. ปรับขนาดอัตโนมัติตามพื้นที่หน้าจอด้วย LayoutBuilder (เส้นผ่านศูนย์กลาง 170 ถึง 240 px)
3. ปรับเปลี่ยนสีเส้นโค้งเรืองแสงอัตโนมัติตามเกณฑ์ปฐพีวิทยา ได้แก่
   - กรดจัด (pH < 5.0) สีแดงเพลิง หรือสีส้มอำพัน
   - เหมาะสมสำหรับทุเรียน (pH 5.5 - 6.5) สีเขียวมรกตนีออน
   - ด่างจัด (pH > 7.5) สีฟ้าคราม
4. มีการ์ดแบบ Frosted Glass พื้นหลังสีกึ่งโปร่งแสง พร้อมขอบเส้นนีออนบางเบา
5. แสดงตัวเลขทศนิยม 2 ตำแหน่งขนาดใหญ่ชัดเจน พร้อมป้ายเปรียบเทียบค่าดิบกับค่าที่ผ่าน AI PINN
6. รองรับการกดแตะเพื่อเปิดหน้าต่างแสดงคำแนะนำอัตราปูนโดโลไมต์ปรับปรุงดิน
\end{lstlisting}
\end{promptbox}

\section{ชุดพรอมพต์แก้ไขปัญหาการเชื่อมต่อ RS485 Half-Duplex บนฮาร์ดแวร์จริง}

\begin{promptbox}[title={พรอมพต์ ก.5 การแก้ปัญหาสายแปลง USB-RS485 มิวต์ตัวรับสัญญาณ (RTS DE/RE)}]
\begin{lstlisting}
คุณคือผู้เชี่ยวชาญด้านฮาร์ดแวร์ไดรเวอร์ USB Serial บนระบบปฏิบัติการแอนดรอยด์
บริบท แอปพลิเคชันเชื่อมต่อสายแปลง USB-to-RS485 (CH340/MAX485) ผ่าน USB OTG สำเร็จ
แต่ไม่ได้รับข้อมูลตอบกลับจากหัววัดดิน (Rx = 0 Bytes ตลอดเวลา)
จงวิเคราะห์สาเหตุและเขียนโค้ดภาษา Dart เพื่อแก้ไขปัญหา ดังนี้
1. อธิบายกลไกการทำงานของสายสัญญาณ RTS (Request to Send) ที่เชื่อมกับขา DE/~RE ของชิป MAX485
2. แก้ไขโค้ดการกำหนดค่าพอร์ต UsbPort ดังนี้
   - ต้องสั่ง port.setRTS(false) เพื่อปลดให้ขา RE เป็น Low เปิดรับสัญญาณจากหัววัด
   - ต้องกำหนด Flow Control เป็น FLOW_CONTROL_OFF
3. เพิ่มระบบ Auto-Baud Rate ตรวจจับและสลับระหว่าง 4800 bps และ 9600 bps อัตโนมัติ
4. เพิ่มระบบเฝ้าระวัง Watchdog Timer ขนาด 3 วินาที กู้คืนการเชื่อมต่อเมื่อไม่มีข้อมูลตอบกลับ
\end{lstlisting}
\end{promptbox}
